\begin{claim}{D.38}{Proved}{A nonnilpotent finite obstruction on the universal \(A_5\) quotient}
Let $R$ be the intersection of the kernels of all continuous
epimorphisms $\widehat F(x,y)\twoheadrightarrow A_5$.
This is an open characteristic subgroup, and the quotient
specified by the nineteen generating-pair representatives in
the accompanying program is
\[
                 Q=F/R\xrightarrow{\ \sim\ } A_5^{19}.
                                                               \tag{D.38a}
\]
Write $X,Y,Z=(XY)^{-1}$ for the images of the three cusps,
and $\theta,\omega$ for the automorphisms of $Q$ induced by
$x\leftrightarrow y$ and $x\mapsto y,\ y\mapsto(xy)^{-1}$.
The image of every full $GT_0$ involution belongs to the
following explicitly computed finite set:
\[
\begin{split}
 \mathcal T=\{t\in Q:\quad&
 t^2=1,\quad tXt^{-1}=X^{-1},\\
 &\omega(t)=Y^{-1}\theta(t)\}.
\end{split}                                                     \tag{D.38b}
\]
Its normalized coordinates are
$\alpha_Q=\operatorname{Inn}(t)$ and
$f_Q=\theta(t)t^{-1}$.
Every tuple in $\mathcal T$ satisfies the finite
normalized relations (I), (II), and
$f_Q\alpha_Q(f_Q)=1$, and defines an actual automorphism
of $Q$ of order two.

The executed exact enumeration gives
\[
\begin{array}{l|r}
\text{finite set or test}&\text{cardinality}\\ \hline
\mathcal T&10\,125\,000\\
\text{tuples surviving the full cusp-triangle obstruction}
       &5\,062\,500\\
\text{tuples conjugate to standard inversion in the finite
pure model (D.38f)}&2\,531\,250\\
\text{triangle survivors outside that finite conjugacy set}
       &2\,531\,250 .
\end{array}                                                     \tag{D.38c}
\]
Thus $5\,062\,500$ explicit finite tuples satisfying the
finite involution equations have no full pure-involution
lift.  The last row is disjoint from the image of the
standard full conjugacy class.  No tuple in that last
row is asserted to have a full lift.
\end{claim}
\begin{proof}
The program enumerates all $60^2$ ordered pairs in $A_5$,
tests generation by closure under multiplication, and
obtains $2\,280$ generating pairs.
Simultaneous conjugation by all $120$ elements of $S_5$
has nineteen free orbits on this set.
These are all $\Aut(A_5)$-orbits: an automorphism
permutes the five Sylow-$2$ subgroups, and equivariance
with the faithful inner $A_5$ action makes that
permutation representation of $\Aut(A_5)$ injective;
conjugation by $S_5$ already supplies its $120$ elements,
as checked in D.31.
Two epimorphisms to $A_5$ have the same kernel exactly
when they differ by an automorphism of $A_5$.
The nineteen kernels are therefore distinct.
The joint quotient map is a subdirect product of
nineteen copies of the nonabelian simple group $A_5$.
The elementary subdirect-product argument shows that
a proper subdirect product of copies of one simple
group identifies two coordinates by an automorphism:
induct on the number of factors and apply the
normal-subgroup description of a direct product of
nonabelian simple groups to the last projection.
Such an identification would equate two of the
nineteen kernels.  This proves surjectivity of the
specified map to $A_5^{19}$ and hence (D.38a).
The kernel is characteristic because every automorphism
of $F$ permutes all epimorphisms to $A_5$.

For each representative pair $(X_i,Y_i)$ the program
also records the ordered triple of conjugacy classes
of $X_i,Y_i,Z_i$ in $A_5$, allowing the simultaneous
exchange of the two classes of elements of order five
by odd conjugation in $S_5$.
All nineteen resulting triples are different.
Every triple contains an element of order five,
and has at least two elements of odd order.
These assertions are checked by complete finite
enumeration, with representatives and component
permutations retained in the output.

By D.34 a full pure involution has multiplier $-1$.
Inversion preserves each conjugacy class of $A_5$:
this is immediate for involutions, and reversal of
a $3$-cycle or $5$-cycle is achieved by an even
permutation.
An automorphism of $A_5^{19}$ permutes its minimal
normal simple factors and acts on each factor by
an element of $\Aut(A_5)=S_5$.
Preservation of the three cusp classes therefore
forces its factor permutation to preserve the
recorded class triple.  The nineteen distinct
triples force the permutation to be the identity.
Since each triple contains an order-five entry,
the component automorphisms must have even parity;
odd conjugation exchanges its two order-five
classes.  Thus the induced automorphism of $Q$
is $\operatorname{Inn}(t)$ for a unique $t\in Q$.
The group $Q$ has trivial center, so its
involutivity gives $t^2=1$.

The exact Harbater--Schneps automorphism identities
from \cite[Propositions 3--4]{R-HS}, with
$m=(-1-1)/2=-1$, descend to $Q$:
\[
 \theta\alpha_Q\theta^{-1}\alpha_Q^{-1}
      =\operatorname{Inn}(f_Q),\qquad
 \omega\alpha_Q\omega^{-1}\alpha_Q^{-1}
      =\operatorname{Inn}(Y^{-1}f_Q).
\]
Center-freeness yields
$f_Q=\theta(t)t^{-1}$ and
$\omega(t)t^{-1}=Y^{-1}f_Q$.
Together with the normalized equation on $X$ these
give precisely (D.38b).

Conversely take $t$ in (D.38b) and define
$f=\theta(t)t^{-1}$.
Then $\theta(t)$ inverts $Y$, so
\[
 tYt^{-1}=f^{-1}Y^{-1}f,\qquad
 \theta(f)f=1,\qquad f\operatorname{Inn}(t)(f)=1.
\]
The second relation in (D.38b) involving $\omega$
can be rewritten $f=Y\omega(t)t^{-1}$.
Apply $\omega$ and $\omega^2$ and use $\omega^3=1$:
\[
 \omega(f)=Z\omega^2(t)\omega(t)^{-1},\qquad
 \omega^2(f)=Xt\omega^2(t)^{-1}.
\]
Substitution now proves the complete finite
three-cycle relation by cancellation,
\[
       \omega^2(f)Z^{-1}\omega(f)Y^{-1}fX^{-1}=1.
                                                               \tag{D.38d}
\]
The tuple $X$ has components of odd order, so
$tXt^{-1}=X^{-1}$ excludes $t=1$ in $Q$.
Thus $\operatorname{Inn}(t)$ has exact order two.
This verifies the assertions for every member
of the finite set, independently of any lifting.

We next impose an obstruction from the full
group.  For each component the extension quotient
is explicitly
\[
 A_5\rtimes_{\operatorname{Inn}(t_i)}C_2
       \xrightarrow{\ \sim\ }A_5\times C_2,\qquad
                (g,s^j)\longmapsto(g t_i^j,s^j).
                                                               \tag{D.38e}
\]
The three vertices of the full cusp triangle in
D.34 are represented by $s$, $xs$, and $yf s$.
Their first coordinates under (D.38e) are
$t_i$, $X_i t_i$, and $Y_i f_i t_i$.
None can be the identity in $A_5$.
Indeed that would make the image of the
corresponding section central in $A_5\times C_2$.
This section class is incident to two different
cusp edges of the full triangle.  At either
edge choose its representative in the cusp
normalizer.  Its image is still central, because
centrality is unchanged by conjugation.
It must therefore both act trivially on the
cusp image and act by inversion there.
Both incident cusp images would have order
at most two.  No cusp image is trivial for
a generating $A_5$ pair.  Two involutions
generate a dihedral group, while any two
entries of the generating triple generate
$A_5$.  This is a contradiction.
Possible fusion of distinct section classes
in the finite quotient does not change this
argument.

For all components together, this obstruction
is exactly
$t_i\notin\{1,X_i\}$ for every $i$.
The third vertex condition follows from these:
$f_i t_i=\theta(t)_i$ and $\theta(X)=Y$,
so $Y_i\theta(t)_i=1$ is another occurrence
of $t_j=X_j$ under the component action of
$\theta$.

Here are the complete finite counts, with the
factorization used to execute them.
The relation $\omega(t)=Y^{-1}\theta(t)$
joins components in twelve independent
sets, using indices $0,\ldots,18$:
\[
 \{0,2\},\{1\},\{3,5\},\{4\},\{6,7\},\{8\},
 \{9,17\},\{10\},\{11,14\},\{12,18\},\{13,15\},\{16\}.
\]
Exhaustive propagation of each initial component
value through its specified $S_5$ conjugation
maps gives respectively
\[
                 3,3,3,3,4,2,5,5,5,5,5,5
\]
assignments.  The equations on a component set
involve no other set, so their Cartesian product
has exactly $10\,125\,000$ elements.
The triangle obstruction changes only the
count for $\{6,7\}$, from four to two.
Its surviving cardinality is therefore
$5\,062\,500$.
All forty-eight component assignments are
executed and checked, including the original
equations (D.38d) and the finite norm; the
Cartesian-product count is exact.

For precision, the finite conjugacy model in
(D.38c) consists of the following actual
automorphisms of $Q$ with a common unit exponent:
\[
\begin{split}
 \mathcal P_Q=\{(\mu,\gamma)\in
      (\Z/30)^\times\times\Aut(Q):\quad&
 \gamma(X)=X^\mu,\\
 &\gamma(Y)\sim_QY^\mu,\quad
 \gamma(Z)\sim_QZ^\mu\}.
\end{split}                                                     \tag{D.38f}
\]
Every normalized full pure automorphism
maps to this set, since $R$ is characteristic.
The recorded class triples again imply that
an element of (D.38f) acts on the $i$th
factor by conjugation by some $g_i\in S_5$,
without permuting factors.
Its parity is even for $\mu\bmod5\in\{1,4\}$
and odd for $\mu\bmod5\in\{2,3\}$.
The remaining condition is exactly
$g_iX_i g_i^{-1}=X_i^\mu$; with that parity,
the $Y_i$ and $Z_i$ class conditions follow.

The standard inversion induces
$\operatorname{Inn}(d)$ on $Q$, where the
program supplies each simultaneous inverter
$d_i$ of $(X_i,Y_i)$.
The conjugacy condition is the finite equation
$t_i=g_i d_i g_i^{-1}$ for every $i$.
The program tests all $120$ choices of $g_i$
for each of the eight units modulo thirty.
There are exactly two distinct compatible
patterns, each with component counts
\[
                  3,3,3,3,1,1,5,5,5,5,5,5.
\]
They are disjoint, so their total cardinality
is $2\cdot1\,265\,625=2\,531\,250$.
This proves the last two rows of (D.38c).
It proves an upper constraint on the image
of the known full conjugacy class, not
surjectivity from that class onto the
finite model.

The complete implementation is
\path{code/part_a/a5_involution_quotient.py};
its output retains the nineteen pairs,
the $\theta,\omega$ component maps, the
standard inversion tuple, every component
count, and an explicit rejected tuple.
No family of lifts through all further
characteristic quotients has been supplied
for a surviving tuple.
\end{proof}
